\documentclass[12pt,a4paper]{article}

% Encoding and fonts
\usepackage[utf8]{inputenc}
\usepackage[T1]{fontenc}
\usepackage{lmodern}
\usepackage{textcomp}

% Page layout
\usepackage[top=1in, bottom=1in, left=1in, right=1in]{geometry}

% Typography
\usepackage{microtype}
\usepackage{parskip}

% Language
\usepackage[english]{babel}

% Headers and footers
\usepackage{fancyhdr}
\pagestyle{fancy}
\fancyhf{}
\fancyhead[L]{\leftmark}
\fancyhead[R]{\thepage}
\fancyfoot[C]{\thepage}
\renewcommand{\headrulewidth}{0.4pt}
\renewcommand{\footrulewidth}{0pt}

% Section styling
\usepackage{titlesec}
\titleformat{\section}{\Large\bfseries}{\thesection}{1em}{}
\titleformat{\subsection}{\large\bfseries}{\thesubsection}{1em}{}
\titleformat{\subsubsection}{\normalsize\bfseries}{\thesubsubsection}{1em}{}
\titlespacing*{\section}{0pt}{2.5ex plus 1ex minus .2ex}{1.5ex plus .2ex}
\titlespacing*{\subsection}{0pt}{2ex plus 1ex minus .2ex}{1ex plus .2ex}
\titlespacing*{\subsubsection}{0pt}{1.5ex plus 1ex minus .2ex}{0.8ex plus .2ex}

% List formatting
\usepackage{enumitem}
\setlist[itemize]{topsep=0.5ex, itemsep=0.25ex, parsep=0pt}
\setlist[enumerate]{topsep=0.5ex, itemsep=0.25ex, parsep=0pt}

% Essential packages
\usepackage{graphicx}
\usepackage[table]{xcolor}
\usepackage{booktabs}
\usepackage{array}
\usepackage{tabularx}
\usepackage{longtable}
\usepackage{amsmath}
\usepackage{amssymb}
\usepackage[normalem]{ulem}
\rowcolors{2}{gray!10}{white}
\renewcommand{\arraystretch}{1.3}

% Cross-references
\usepackage{hyperref}
\usepackage{cleveref}

% Custom commands
\newcommand{\pilcrow}{\S}

% Hyperref setup
\hypersetup{
    colorlinks=true,
    linkcolor=blue,
    filecolor=magenta,
    urlcolor=cyan,
}

% Code listings
\usepackage{listings}
\definecolor{codebg}{RGB}{248,248,248}
\definecolor{codeframe}{RGB}{200,200,200}
\definecolor{codekeyword}{RGB}{0,0,180}
\definecolor{codecomment}{RGB}{0,128,0}
\definecolor{codestring}{RGB}{163,21,21}
\lstset{
    basicstyle=\ttfamily\small,
    breaklines=true,
    breakatwhitespace=false,
    frame=single,
    framerule=0.5pt,
    rulecolor=\color{codeframe},
    backgroundcolor=\color{codebg},
    keepspaces=true,
    numbers=left,
    numberstyle=\tiny\color{gray},
    numbersep=8pt,
    showstringspaces=false,
    tabsize=4,
    xleftmargin=1em,
    xrightmargin=1em,
    aboveskip=1em,
    belowskip=1em,
    keywordstyle=\color{codekeyword}\bfseries,
    commentstyle=\color{codecomment}\itshape,
    stringstyle=\color{codestring},
    literate={_}{{\textunderscore}}1
             {-}{{\textminus}}1
             {<}{{\textless}}1
             {>}{{\textgreater}}1
             {~}{{\textasciitilde}}1
             {^}{{\textasciicircum}}1
}

% YAML language definition for listings
\lstdefinelanguage{YAML}{
    keywords={true,false,null,y,n},
    keywordstyle=\color{blue}\bfseries,
    sensitive=false,
    comment=[l]{\#},
    morecomment=[s]{/*}{*/},
    commentstyle=\color{gray}\ttfamily,
    stringstyle=\color{red}\ttfamily,
    morestring=[b]'',
    morestring=[b]'',
}

% TOML language definition for listings
\lstdefinelanguage{TOML}{
    keywords={true,false},
    keywordstyle=\color{blue}\bfseries,
    sensitive=true,
    comment=[l]{\#},
    commentstyle=\color{gray}\ttfamily,
    stringstyle=\color{red}\ttfamily,
    morestring=[b]'',
    morestring=[b]'',
}

% Dockerfile language definition for listings
\lstdefinelanguage{Dockerfile}{
    keywords={FROM,RUN,CMD,LABEL,EXPOSE,ENV,ADD,COPY,ENTRYPOINT,VOLUME,USER,WORKDIR,ARG,ONBUILD,STOPSIGNAL,HEALTHCHECK,SHELL},
    keywordstyle=\color{blue}\bfseries,
    sensitive=false,
    comment=[l]{\#},
    commentstyle=\color{gray}\ttfamily,
    stringstyle=\color{red}\ttfamily,
    morestring=[b]'',
    morestring=[b]'',
}

% JSON language definition for listings
\lstdefinelanguage{JSON}{
    keywords={true,false,null},
    keywordstyle=\color{blue}\bfseries,
    sensitive=true,
    commentstyle=\color{gray}\ttfamily,
    stringstyle=\color{red}\ttfamily,
    morestring=[b]'',
}

% Code listing float environment
\usepackage{float}
\newfloat{listing}{htbp}{lol}[section]
\floatname{listing}{Listing}

% Admonition boxes
\usepackage{tcolorbox}
\tcbuselibrary{skins,breakable}

% Admonition style definitions
\newtcolorbox{admonitionbox}[2][]{
    enhanced,
    breakable,
    colback=#2!5,
    colframe=#2!80!black,
    fonttitle=\bfseries,
    title=#1,
    left=2mm,
    right=2mm,
}

% Synthon tuple boxes
\newtcolorbox{synthonbox}{
    enhanced,
    colback=blue!5,
    colframe=blue!75!black,
    fontupper=\large,
    center upper,
    boxrule=1pt,
    arc=4pt,
}

\begin{document}

\title{The Probability That Two Random Elements Generate the Symmetric Group}
\date{\today}

\maketitle

\subsection{The Exact Formula (Hall''s Identity)}

Let $G = S_n$ be the symmetric group on $n$ letters, with $|G| = n!$. Let $\mu$ denote the
Möbius function of the subgroup lattice of $G$. The probability that two elements chosen
uniformly and independently from $S_n$ generate $S_n$ is:

\[P_n = \frac{1}{(n!)^2}\sum_{H \leq S_n} \mu(H)\, |H|^2, \tag{1}\]

where the sum runs over all subgroups $H$ of $S_n$, and $\mu: \text{Sub}(G) \to \mathbb{Z}$
is defined recursively by $\mu(G) = 1$ and

\[\mu(H) = -\sum_{H < K \leq G} \mu(K). \tag{2}\]

This is Hall''s formula (1936) for any finite group $G$.

\subsection{Verification for Small $n$}

\subsubsection{$n = 3$}

The subgroup lattice of $S_3$ (order 6):

\textbar{} Subgroup $H$ \textbar{} $|H|$ \textbar{} $\mu(H)$ \textbar{} $\mu(H)|H|^2$ \textbar{}
\textbar{}---\textbar{}---\textbar{}---\textbar{}---\textbar{}
\textbar{} $S_3$ \textbar{} 6 \textbar{} 1 \textbar{} 36 \textbar{}
\textbar{} $A_3 \cong C_3$ \textbar{} 3 \textbar{} $-1$ \textbar{} $-9$ \textbar{}
\textbar{} $\langle(12)\rangle \cong C_2$ \textbar{} 2 \textbar{} $-1$ \textbar{} $-4$ \textbar{}
\textbar{} $\langle(13)\rangle \cong C_2$ \textbar{} 2 \textbar{} $-1$ \textbar{} $-4$ \textbar{}
\textbar{} $\langle(23)\rangle \cong C_2$ \textbar{} 2 \textbar{} $-1$ \textbar{} $-4$ \textbar{}
\textbar{} $\{e\}$ \textbar{} 1 \textbar{} 3 \textbar{} 3 \textbar{}

\[\sum_H \mu(H)|H|^2 = 36 - 9 - 12 + 3 = 18, \quad P_3 = \frac{18}{36} = \frac{1}{2}.\]

Indeed, in $S_3$: an element is even with probability $1/2$. Two random elements generate
$S_3$ iff they are not both in $A_3$ (probability $3/4$) and not both in the same order-2
subgroup. Checking: pairs of odd permutations with at least one transposition always generate,
as do mixed parity pairs with appropriate support.

Verification: the $18$ generating pairs out of $36$ are confirmed by brute force.

\subsubsection{$n = 4$}

$|S_4| = 24$, $24^2 = 576$. The subgroup lattice is richer (30 subgroups up to
isomorphism, with many conjugacy classes). The sum yields $216$ generating pairs:

\[P_4 = \frac{216}{576} = \frac{3}{8} = 0.375.\]

\subsubsection{$n = 5$}

$|S_5| = 120$, $120^2 = 14400$. The sum yields $6840$ generating pairs:

\[P_5 = \frac{6840}{14400} = \frac{19}{40} = 0.475.\]

\subsection{Organizing by Conjugacy Classes of Subgroups}

To compute (1) efficiently, group the sum by conjugacy classes of maximal subgroups.
The maximal subgroups of $S_n$ fall into three types (the O''Nan--Scott theorem):

\subsubsection{Type 1: Intransitive subgroups}

\[M \cong S_k \times S_{n-k}, \quad 1 \leq k \leq \left\lfloor\frac{n}{2}\right\rfloor.\]


Number of conjugates: $\binom{n}{k}$ (for $k < n/2$) or $\frac{1}{2}\binom{n}{n/2}$ (for $k = n/2$).
Each has order $k! \,(n-k)!$.

\subsubsection{Type 2: Imprimitive subgroups}

\[M \cong S_k \wr S_m \quad \text{where } n = k \cdot m, \quad 1 < k < n.\]


These are stabilizers of a nontrivial partition of $\{1,\dots,n\}$ into $m$ blocks of size $k$.

\subsubsection{Type 3: Primitive subgroups}

Maximal subgroups acting primitively on $\{1,\dots,n\}$ (classified using the
Classification of Finite Simple Groups). These include $A_n$, and various affine
and almost-simple groups.

\subsubsection{The dominant term: the alternating group $A_n$}

The unique index-2 subgroup $A_n$ has $\mu(A_n) = -1$ (it sits immediately below
$S_n$ in the lattice). Its contribution to (1) is:

\[\frac{\mu(A_n) \cdot |A_n|^2}{(n!)^2} = \frac{-1 \cdot (n!/2)^2}{(n!)^2} = -\frac{1}{4}.\]

This term alone contributes $-\tfrac{1}{4}$ to $P_n$.

\subsection{The Asymptotic Formula}

Dixon (1969) proved the landmark result:

\[P(S_n \text{ or } A_n \text{ is generated}) \to 1 \quad \text{as } n \to \infty.\]

Since $P(\text{both elements land in } A_n) = \left(\frac{1}{2}\right)^2 = \frac{1}{4}$,
and conditional on avoiding $A_n$ the pair generates $S_n$ with probability tending
to 1, we obtain:

\[\boxed{\lim_{n \to \infty} P_n = \frac{3}{4}.}\]

Bovey (1980), refined by Maróti (2011), gave the precise asymptotic expansion:

\[P_n = \frac{3}{4} - \frac{1}{n!} - \frac{1}{2\binom{n}{n/2}} + O\left(n^{-3/2}\log n\right),\]

where the second term $\frac{1}{2\binom{n}{n/2}}$ is omitted when $n$ is odd.
Using Stirling''s approximation, we have $\binom{n}{n/2} \sim \frac{2^n}{\sqrt{\pi n/2}}$, so

\[P_n = \frac{3}{4} - \frac{1}{n!} - \frac{\sqrt{\pi n}}{2 \cdot 4^n} + O\left(n^{-3/2}\log n\right). \tag{3}\]

The convergence to $\tfrac{3}{4}$ is extremely rapid --- already by $n = 5$,
$P_5 = 0.475$, and the error shrinks exponentially thereafter.

\subsection{Organizing by Maximal Subgroups}

Applying Hall identity (1) directly requires the full subgroup lattice, which is
enormous for large $n$. A more practical approach groups the sum by maximal subgroups.

By Mobius inversion on the subgroup lattice, we have the equivalent formula:

\[P_n = \frac{1}{(n!)^2} \sum_{\substack{M \text{ maximal} \\ M \neq S_n,\, M \neq A_n}} \mu(M)\, |M|^2 + \frac{1}{(n!)^2} \sum_{\substack{H \text{ non-maximal} \\ H \neq \{e\}}} \mu(H)\, |H|^2.\]

The dominant correction comes from the intransitive maximal subgroups
$S_k \times S_{n-k}$, each of proportion $1/\binom{n}{k}$ in $S_n$. Their total
contribution to the sum is $O(1/n)$ and vanishes as $n \to \infty$.

The primitive maximal subgroups other than $A_n$ contribute $O(n^{-3/2}\log n)$
(Dixon, 1969).

\subsection{Proof Sketch of $\lim_{n \to \infty} P_n = \frac{3}{4}$}

\textbf{Step 1.} By the Classification of Finite Simple Groups and Jordan''s theorem, a
transitive subgroup of $S_n$ containing a transposition is all of $S_n$. A
transitive subgroup containing a 3-cycle and acting 2-transitively is either
$A_n$ or $S_n$.

\textbf{Step 2.} The probability that two random permutations are both even is $1/4$. In
this case they can only generate $A_n$ or a proper subgroup.

\textbf{Step 3.} The probability that two random permutations both lie in a fixed
intransitive maximal subgroup $S_k \times S_{n-k}$ is $\binom{n}{k}^{-2}$. Summing
over all such subgroups gives $\sum_{k=1}^{\lfloor n/2 \rfloor}\binom{n}{k}^{-1} = O(1/n)$.

\textbf{Step 4.} Dixon showed that the probability that two random permutations generate
a transitive subgroup other than $A_n$ or $S_n$ is $O(n^{-3/2}\log n)$.

\textbf{Step 5.} Combining:


\[P_n = 1 - \frac{1}{4} - O(1/n) - O(n^{-3/2}\log n) = \frac{3}{4} + o(1).\]

\subsection{Generalization to $d$ Generators}

For $d$ random generators, the same analysis gives:

\[P(S_n \text{ is generated by } d \text{ elements}) \to 1 - 2^{-d} \quad \text{as } n \to \infty.\]

The $2^{-d}$ comes from the probability that all $d$ elements land in $A_n$.

\begin{center}
\rule{0.5\textwidth}{0.4pt}
\end{center}


\end{document}
